\documentclass[tikz,border=8pt]{standalone}
\usepackage{ctex}
\usepackage{amsmath}
\usetikzlibrary{calc, arrows.meta}

\begin{document}
\begin{tikzpicture}[scale=1.0, thick]

% 定点问题套路示意：含参直线族 y = k(x-1)+2 过定点 (1,2)

% 坐标轴
\draw[->, gray, thin] (-2.0,0) -- (3.8,0) node[right, font=\footnotesize] {$x$};
\draw[->, gray, thin] (0,-1.5) -- (0,4.5) node[above, font=\footnotesize] {$y$};
\node[font=\footnotesize, below left] at (0,0) {$O$};

% 定点 (1,2)
\coordinate (P) at (1,2);

% 五条不同斜率的直线 y = k(x-1)+2，k 分别为 -2, -0.8, 0, 1.5, 3.5
% 每条线画在合理范围内 x ∈ [-1.5, 3.5]
\foreach \k/\col in {
    -2.0/blue!35,
    -0.7/blue!50,
     0.0/blue!65,
     1.5/blue!80,
     3.5/blue!95
}{
  \draw[\col, thin]
    ({-1.5}, {\k*(-1.5-1)+2}) -- ({3.2}, {\k*(3.2-1)+2});
}

% 斜率标注（放在右端或外侧，避免重叠）
\node[font=\tiny, blue!35, right] at (3.2, {-2.0*(3.2-1)+2})  {$k=-2$};
\node[font=\tiny, blue!50, right] at (3.2, {-0.7*(3.2-1)+2})  {$k=-0.7$};
\node[font=\tiny, blue!65, right] at (3.2, {0.0*(3.2-1)+2})   {$k=0$};
\node[font=\tiny, blue!80, right] at (3.2, {1.5*(3.2-1)+2})   {$k=1.5$};
\node[font=\tiny, blue!95, right] at (3.2, {3.5*(3.2-1)+2})   {$k=3.5$};

% 定点 (1,2) 高亮红色
\fill[red] (P) circle [radius=3.5pt];
\node[red, above right, font=\normalsize\bfseries] at (P) {定点 $(1,\,2)$};

% 定点坐标参考线（淡灰虚线）
\draw[gray!50, dashed, very thin] (1,0) node[below, font=\scriptsize, gray] {$1$} -- (1,2);
\draw[gray!50, dashed, very thin] (0,2) node[left, font=\scriptsize, gray] {$2$} -- (1,2);

% 套路说明框（左下角）
\node[draw=orange!60, fill=orange!6, rounded corners=5pt,
      font=\small, align=left, inner sep=8pt, text width=5.5cm]
  at (-0.2,-1.0)
  {\textbf{套路：}\\
   $y=k(x-1)+2$\\
   $\Rightarrow\ k(x-1)+(2-y)=0$\\[2pt]
   令 $x-1=0$ 且 $y-2=0$\\
   $\Rightarrow$ 定点 $(1,2)$（与 $k$ 无关）};

% 标题
\node[font=\normalsize\bfseries] at (1.5,4.1)
  {含参直线族的定点};

\end{tikzpicture}
\end{document}
